% Unit 054 terjemahan Bahasa Indonesia dari chapter4.tex baris 723--986.
% Sumber: Methods of Algebra, Volume 2, commit
% 9a5803ff77dd3257484cb177f851a73770a59dd3 (CC BY 4.0).
% stable-unit-id: o014.aljabr2.chapter4.derived-categories
% Cakupan: Bagian 4.4 mengenai konstruksi dan sifat dasar kategori turunan;
% berhenti sebelum Bagian 4.5 pada baris 987.

% segment-id: o014.li2.chapter4.derived-categories.h001
\section{Kategori Turunan}\label{sec:derived-cat}
\hypertarget{o014.aljabr2.chapter4.derived-categories}{ }

% segment-id: o014.li2.chapter4.derived-categories.p001
Paruh pertama bagian ini membahas kategori aditif $\mathcal{A}$; ketika
kategori turunan didefinisikan kemudian, $\mathcal{A}$ akan disyaratkan
sebagai kategori abelian. Contoh~\ref{eg:cplx-as-translation-cat} telah
menunjukkan bahwa $\cate{C}(\mathcal{A})$ dan $\cate{K}(\mathcal{A})$
merupakan kategori aditif dengan pergeseran; funktor pergeseran keduanya
dituliskan sebagai $X \mapsto X[1]$.

% segment-id: o014.li2.chapter4.derived-categories.p002
Tinjau morfisme $f: X \to Y$ dalam $\cate{C}(\mathcal{A})$ beserta kerucut
pemetaan $\Cone(f)$. Definisi~\ref{def:cone-triangle} kemudian memberikan
segitiga berikut dalam $\cate{C}(\mathcal{A})$:
% segment-id: o014.li2.chapter4.derived-categories.g001
\begin{equation}\label{eqn:Cone-triangle} \begin{tikzcd}
	X \arrow[r, "f"] & Y \arrow[r, "{\alpha(f)}"] & \Cone(f) \arrow[r, "\beta(f)"] & X[1]
\end{tikzcd}\end{equation}
Dengan mengambil kuosien, diperoleh segitiga dalam $\cate{K}(\mathcal{A})$.

% Below: [KS06, Theorem 11.2.6]
% segment-id: o014.li2.chapter4.derived-categories.th001
\begin{teorema}\label{prop:cplx-triangulated}
	Pandang $\cate{K}(\mathcal{A})$ sebagai kategori aditif dengan pergeseran
	$T: X \mapsto X[1]$. Sebut suatu segitiga $X \to Y \to Z \xrightarrow{+1}$
	sebagai segitiga terbedakan dalam $\cate{K}(\mathcal{A})$ apabila segitiga
	tersebut isomorfik dengan segitiga berbentuk \eqref{eqn:Cone-triangle}.
	Segitiga-segitiga terbedakan ini membentuk himpunan $\mathcal{H}$ yang
	menjadikan data $(\cate{K}(\mathcal{A}), T, \mathcal{H})$ suatu kategori
	bertriangulasi.
\end{teorema}

% segment-id: o014.li2.chapter4.derived-categories.pf001
\begin{bukti}
	Kita verifikasi aksioma-aksioma dalam
	Definisi~\ref{def:triangulated-cat} satu per satu. (TR0) dan (TR2) jelas
	berdasarkan definisi, sedangkan (TR3) mengikuti dari diagram komutatif dalam
	$\cate{K}(\mathcal{A})$ yang diberikan oleh
	Lema~\ref{prop:cone-alpha}:
% segment-id: o014.li2.chapter4.derived-categories.g002
	\[\begin{tikzcd}
		Y \arrow[r, "{\alpha(f)}"] \arrow[d, "\identity"'] & \Cone(f) \arrow[r, "{\beta(f)}"] \arrow[d, "\identity"] & X[1] \arrow[r, "{-f[1]}"] \arrow[d, "\sim" sloped] & Y[1] \arrow[d, "\identity"] \\
		Y \arrow[r, "{\alpha(f)}"'] & \Cone(f) \arrow[r, "{\alpha(\alpha(f))}"'] & \Cone(\alpha(f)) \arrow[r, "{\beta(\alpha(f))}"'] & Y[1] .
	\end{tikzcd}\]

	Untuk (TR1), bagi $X \in \Obj(\cate{K}(\mathcal{A}))$ tinjau kerucut
	pemetaan dari morfisme nol $0 \to X \to X \xrightarrow{+1}$; rotasi dengan
	(TR3) memberikan segitiga terbedakan
	$X \xrightarrow{\identity_X} X \to 0 \xrightarrow{+1}$.

	Untuk (TR4), tinjau diagram komutatif berikut dalam
	$\cate{K}(\mathcal{A})$ (bagian bergaris penuh):
% segment-id: o014.li2.chapter4.derived-categories.g003
	\begin{equation}\label{eqn:cplx-triangulated-aux-0}\begin{tikzcd}
		X \arrow[d, "\alpha"'] \arrow[r, "f"] & Y \arrow[r, "\alpha(f)"] \arrow[d, "\beta"] & \Cone(f) \arrow[r, "\beta(f)"] \arrow[dashed, d, "\gamma"] & TX \arrow[dashed, d, "T\alpha"] \\
		X' \arrow[r, "{f'}"'] & Y' \arrow[r, "{\alpha(f')}"'] & \Cone(f') \arrow[r, "{\beta(f')}"'] & T X' .
	\end{tikzcd}\end{equation}
	Komutativitas persegi kiri ekuivalen dengan keberadaan suatu keluarga
	morfisme dalam $\mathcal{A}$, yaitu $h^n: X^n \to (Y')^{n-1}$, sedemikian
	rupa sehingga
% segment-id: o014.li2.chapter4.derived-categories.q001
	\[ \beta^n f^n - (f')^n \alpha^n = h^{n+1} d_X^n + d_{Y'}^{n-1} h^n , \quad n \in \Z . \]
	Dengan notasi matriks, untuk setiap $n \in \Z$ definisikan morfisme
% segment-id: o014.li2.chapter4.derived-categories.q002
	\[ \gamma^n = \begin{pmatrix} \alpha^{n+1} & 0 \\ h^{n+1} & \beta^n \end{pmatrix} : X^{n+1} \oplus Y^n \to (X')^{n+1} \oplus (Y')^n . \]
	Periksa bahwa
% segment-id: o014.li2.chapter4.derived-categories.q003
	\[ \begin{pmatrix} \alpha^{n+1} & 0 \\ h^{n+1} & \beta^n \end{pmatrix} \begin{pmatrix} -d_X^n & 0 \\ f^n & d_Y^{n-1} \end{pmatrix}
	= \begin{pmatrix} -d_{X'}^n & 0 \\ (f')^n & d_{Y'}^{n-1} \end{pmatrix} \begin{pmatrix} \alpha^n & 0 \\ h^n & \beta^{n-1} \end{pmatrix}, \]
	sehingga semua matriks tersebut menentukan morfisme
	$\gamma: \Cone(f) \to \Cone(f')$. Jelas bahwa dua persegi di sebelah kanan
	dalam \eqref{eqn:cplx-triangulated-aux-0} komutatif di
	$\cate{C}(\mathcal{A})$. Perhatikan bahwa homotopi $(h^n)_{n \in \Z}$ dapat
	dipilih dengan berbagai cara, sehingga $\gamma$ tidak unik.

	Tersisa verifikasi (TR5). Dalam pernyataan aksioma tersebut, tanpa mengurangi
	keumuman ambil $Z' = \Cone(f)$, $X' = \Cone(g)$, dan
	$Y' = \Cone(gf)$, sedangkan $h$, $k$, dan $m$ berturut-turut adalah
	$\alpha(f)$, $\alpha(g)$, dan $\alpha(gf)$; morfisme-morfisme berderajat
	$+1$ yang bersesuaian berbentuk $\beta(\cdot)$. Funktorialitas kerucut
	pemetaan (Proposisi~\ref{prop:cone-functorial}), atau pemeriksaan langsung,
	memberikan morfisme-morfisme
% segment-id: o014.li2.chapter4.derived-categories.g004
	\[\begin{tikzcd}[column sep=large, row sep=small]
		\Cone(f) \arrow[r, "{u := (\identity_{X[1]}, g)}" inner sep=0.8em] & \Cone(gf) \arrow[r, "{v := (f[1], \identity_Z)}" inner sep=0.8em] & \Cone(g) \\
		Z' \arrow[equal, u] & Y' \arrow[equal, u] & X' \arrow[equal, u] .
	\end{tikzcd}\]
	Berdasarkan funktorialitas kerucut pemetaan, semua persegi dalam diagram
	(TR5) yang tidak melibatkan $w$ menjadi komutatif. Agar bagian yang tersisa,
	yakni persegi kanan bawah dalam (TR5), juga komutatif, morfisme
	$w: X' \to TZ'$ harus dan hanya dapat diambil sebagai komposisi
% segment-id: o014.li2.chapter4.derived-categories.q004
	\[ X' \xrightarrow{\beta(g)} TY \xrightarrow{T(\alpha(f))} T Z' . \]

	Dengan demikian, tinggal dibuktikan bahwa
	$Z' \xrightarrow{u} Y' \xrightarrow{v} X' \xrightarrow{w} TZ'$ merupakan
	segitiga terbedakan. Perhatikan bahwa untuk setiap $n \in \Z$,
% segment-id: o014.li2.chapter4.derived-categories.q005
	\begin{align*}
		\Cone(u)^n & = \Cone(f)^{n+1} \oplus \Cone(gf)^n = X^{n+2} \oplus Y^{n+1} \oplus X^{n+1} \oplus Z^n , \\
		(X')^n & = \Cone(g)^n = Y^{n+1} \oplus Z^n .
	\end{align*}
	Dengan notasi matriks, definisikan
% segment-id: o014.li2.chapter4.derived-categories.q006
	\begin{align*}
		\varphi^n & := \begin{pmatrix} 0 & \identity_{Y^{n+1}} & f^{n+1} & 0 \\ 0 & 0 & 0 & \identity_{Z^n} \end{pmatrix} : \Cone(u)^n \to (X')^n , \\
		\psi^n & := \begin{pmatrix} 0 & 0 \\ \identity_{Y^{n+1}} & 0 \\ 0 & 0 \\ 0 & \identity_{Z^n} \end{pmatrix} : (X')^n \to \Cone(u)^n .
	\end{align*}
	Pemeriksaan langsung menunjukkan bahwa keduanya memberikan morfisme
	$\varphi: \Cone(u) \leftrightarrows X' : \psi$ dalam
	$\cate{C}(\mathcal{A})$ yang memenuhi
	$\varphi \circ \psi = \identity_{X'}$, serta
	$\varphi \circ \alpha(u) = v$ dan $\beta(u) \circ \psi = w$.
	Selanjutnya, tinjau diagram
% segment-id: o014.li2.chapter4.derived-categories.g005
	\[\begin{tikzcd}
		Y' \arrow[r, "v"] \arrow[d, "\identity_{Y'}"'] & X' \arrow[r, "w"] \arrow[shift left, d, "\psi"] & TZ' \arrow[d, "\identity_{TZ'}"] \\
		Y' \arrow[r, "\alpha(u)"'] & \Cone(u) \arrow[r, "\beta(u)"'] \arrow[shift left, u, "\varphi"] & TZ'
	\end{tikzcd}\]
	Cukup ditunjukkan bahwa diagram tersebut komutatif dalam
	$\cate{K}(\mathcal{A})$ serta bahwa $\psi$ dan $\varphi$ saling invers dalam
	$\cate{K}(\mathcal{A})$; satu-satunya hal yang masih perlu dibuktikan ialah
	bahwa $\psi \circ \varphi$ homotopik dengan
	$\identity_{\Cone(u)}$. Ambil $h^n : \Cone(u)^n \to \Cone(u)^{n-1}$ sebagai
	komposisi langsung
% segment-id: o014.li2.chapter4.derived-categories.q007
	\[ X^{n+2} \oplus Y^{n+1} \oplus X^{n+1} \oplus Z^n \twoheadrightarrow X^{n+1} \hookrightarrow X^{n+1} \oplus Y^n \oplus X^n \oplus Z^{n-1}. \]
	Verifikasi rutin memberikan
	$\identity_{\Cone(u)^n} - \psi^n \varphi^n = h^{n+1} d_{\Cone(u)}^n + d_{\Cone(u)}^{n-1} h^n$.
	Pernyataan pun terbukti.
\end{bukti}

% segment-id: o014.li2.chapter4.derived-categories.p003
Ekuivalensi homotopi sangat diperlukan dalam bukti tersebut. Jika kita bekerja
pada tingkat $\cate{C}(\mathcal{A})$, kerucut pemetaan tidak dapat memberikan
suatu kategori bertriangulasi.

% segment-id: o014.li2.chapter4.derived-categories.pf002
\begin{bukti}[Bukti Teorema~\ref{prop:resolution-homotopy} dan~\ref{prop:K-resolution-homotopy}]
	Sekarang kita dapat melengkapi bagian bukti
	Teorema~\ref{prop:resolution-homotopy} yang sebelumnya tertunda. Berdasarkan
	dualitas, cukup ditinjau keadaan $Y \xleftarrow{\alpha} X \xrightarrow{\gamma} I$,
	dengan $\alpha$ suatu kuasi-isomorfisme dan
	$I \in \Obj(\cate{C}^+(\mathcal{A}))$ terdiri atas objek-objek injektif.
	Funktor kohomologis $\Hom_{\cate{K}(\mathcal{A})}(\cdot, I)$
	(Proposisi~\ref{prop:cohomological-Hom-functor}) memberikan barisan eksak
	panjang
% segment-id: o014.li2.chapter4.derived-categories.q008
	\begin{multline*}
		\Hom_{\cate{K}(\mathcal{A})}\left( \Cone(\alpha), I \right)  \to \Hom_{\cate{K}(\mathcal{A})}(Y, I) \\
		\xrightarrow{\alpha^*} \Hom_{\cate{K}(\mathcal{A})}(X, I) \to \Hom_{\cate{K}(\mathcal{A})}(\Cone(\alpha)[-1], I);
	\end{multline*}
	Akan tetapi, $\Cone(\alpha)$ asiklik
	(Korolari~\ref{prop:quism-cone}), sehingga
	Lema~\ref{prop:resolution-homotopy-aux} menyiratkan bahwa $\alpha^*$ adalah
	isomorfisme.

	Adapun untuk Teorema~\ref{prop:K-resolution-homotopy} mengenai kompleks
	K-injektif atau K-projektif, argumennya sepenuhnya serupa.
\end{bukti}

% segment-id: o014.li2.chapter4.derived-categories.co001
\begin{korolari}
	Untuk $\star \in \{+, -, \bdd\}$, kategori
	$\cate{K}^\star(\mathcal{A})$ yang diperkenalkan dalam
	Definisi~\ref{def:cplx-homotopy-cat-variant} merupakan subkategori
	bertriangulasi dari $\cate{K}(\mathcal{A})$
	(Definisi--Proposisi~\ref{def:triangulated-subcat}).
\end{korolari}

% segment-id: o014.li2.chapter4.derived-categories.pf003
\begin{bukti}
	Jelas bahwa $\cate{K}^\star(\mathcal{A})$ tertutup terhadap funktor
	pergeseran. Selanjutnya, bagi kerucut pemetaan
	$X \xrightarrow{f} Y \to \Cone(f) \xrightarrow{+1}$, jelas bahwa $X,Y$
	termasuk dalam $\cate{C}^\star(\mathcal{A})$ jika dan hanya jika
	$\Cone(f) \in \cate{C}^\star(\mathcal{A})$. Hal ini membuktikan bahwa syarat
	(ii) dalam Definisi--Proposisi~\ref{def:triangulated-subcat} terpenuhi.
\end{bukti}

% segment-id: o014.li2.chapter4.derived-categories.pr001
\begin{proposisi}\label{prop:functor-induces-triangulated}
	\index[sym1]{KF}
	Misalkan $F: \mathcal{A}_1 \to \mathcal{A}_2$ suatu funktor aditif di antara
	kategori-kategori aditif. Maka
	$\cate{K}F: \cate{K}(\mathcal{A}_1) \to \cate{K}(\mathcal{A}_2)$ merupakan
	funktor bertriangulasi. Pernyataan yang sama berlaku dengan
	$\cate{K}^\star(\cdot)$ menggantikan $\cate{K}(\cdot)$, dengan
	$\star \in \{+, -, \bdd\}$.
\end{proposisi}

% segment-id: o014.li2.chapter4.derived-categories.pf004
\begin{bukti}
	Pernyataan ini mengikuti dari Proposisi~\ref{prop:CA-functoriality}
	(mengenai pergeseran) dan Proposisi~\ref{prop:Cone-A}
	(mengenai kerucut pemetaan).
\end{bukti}

% segment-id: o014.li2.chapter4.derived-categories.co002
\begin{korolari}\label{prop:subcat-K-triangulated}
	Misalkan $\mathcal{A}'$ suatu subkategori penuh aditif dari $\mathcal{A}$.
	Maka $\cate{K}(\mathcal{A}')$ membenam sebagai subkategori bertriangulasi
	dari $\cate{K}(\mathcal{A})$. Pernyataan yang sama berlaku dengan
	$\cate{K}^\star(\cdot)$ menggantikan $\cate{K}(\cdot)$, dengan
	$\star \in \{+, -, \bdd\}$.
\end{korolari}

% segment-id: o014.li2.chapter4.derived-categories.pr002
\begin{proposisi}
	Misalkan $\mathcal{A}$ suatu kategori abelian. Untuk setiap $n \in \Z$,
	funktor $\Hm^n: \cate{K}(\mathcal{A}) \to \mathcal{A}$ merupakan funktor
	kohomologis dalam pengertian Definisi~\ref{def:cohomological-functor}.
	Pernyataan yang sama berlaku dengan $\cate{K}^\star(\mathcal{A})$
	menggantikan $\cate{K}(\mathcal{A})$, dengan
	$\star \in \{+, -, \bdd\}$.
\end{proposisi}

% segment-id: o014.li2.chapter4.derived-categories.pf005
\begin{bukti}
	Tinjau segitiga terbedakan dalam $\cate{K}(\mathcal{A})$ yang diberikan oleh
	kerucut pemetaan,
	$X \xrightarrow{f} Y \xrightarrow{\alpha(f)} \Cone(f) \xrightarrow{+1}$.
	Berdasarkan barisan eksak panjang
	\eqref{eqn:cone-long-exact-sequence} dan
	Korolari~\ref{prop:cone-connecting} (inti dari
	\S\ref{sec:cone-vs-long-exact-sequence}), barisan
% segment-id: o014.li2.chapter4.derived-categories.q009
	\[ \Hm^n(X) \xrightarrow{\Hm^n(f)} \Hm^n(Y) \xrightarrow{\Hm^n(\alpha(f))} \Hm^n(\Cone(f)) \]
	memang eksak.
\end{bukti}

% segment-id: o014.li2.chapter4.derived-categories.p004
Definisi kategori turunan akan melibatkan lokalisasi Verdier. Dengan notasi
Contoh~\ref{eg:cohomological-functor-multiplicative}, funktor kohomologis
$\Hm^0: \cate{K}(\mathcal{A}) \to \mathcal{A}$ menentukan sistem
multiplikatif yang kompatibel dengan triangulasi
\index[sym1]{qis@$\mathrm{qis}$}
% segment-id: o014.li2.chapter4.derived-categories.q010
\[ \mathrm{qis} := S_{\Hm^0} \subset \Mor(\cate{K}(\mathcal{A})), \]
	yang unsur-unsurnya tepat merupakan kuasi-isomorfisme. Sudut pandang lain
untuk memahami $\mathrm{qis}$ ialah meninjau subkategori bertriangulasi jenuh
$\mathcal{N} := \mathcal{N}_{\Hm^0}$ dari $\cate{K}(\mathcal{A})$ yang terdiri
atas kompleks-kompleks asiklik. Proposisi~\ref{prop:cohomological-functor-multiplicative}
menyiratkan $S\mathcal{N} = \mathrm{qis}$.

% segment-id: o014.li2.chapter4.derived-categories.p005
Untuk $\star \in \{+, -, \bdd\}$, secara serupa kita dapat mendefinisikan
sistem multiplikatif yang kompatibel dengan triangulasi
$\mathrm{qis}^\star := \mathrm{qis} \cap \Mor(\cate{K}^\star(\mathcal{A}))$
dan $\mathcal{N}^\star := \mathcal{N} \cap \cate{K}^\star(\mathcal{A})$;
keduanya tetap memenuhi
$S\mathcal{N}^\star = \mathrm{qis}^\star$.

% segment-id: o014.li2.chapter4.derived-categories.d001
\begin{definisi}\label{def:derived-cat}
	\index{kategori turunan (derived category)}
	\index[sym1]{DA@$\cate{D}(\mathcal{A})$, $\cate{D}^+(\mathcal{A})$, $\cate{D}^-(\mathcal{A})$, $\cate{D}^{\bdd}(\mathcal{A})$}
	\emph{Kategori turunan} dari suatu kategori abelian $\mathcal{A}$
	didefinisikan sebagai kategori bertriangulasi
% segment-id: o014.li2.chapter4.derived-categories.q011
	\begin{align*}
		\cate{D}(\mathcal{A}) & := \cate{K}(\mathcal{A})\left[ \mathrm{qis}^{-1} \right] \\
		& = \cate{K}(\mathcal{A})/\mathcal{N}.
	\end{align*}
	Demikian pula, untuk $\star \in \{+, -, \bdd\}$ dapat didefinisikan kategori
	bertriangulasi
% segment-id: o014.li2.chapter4.derived-categories.q012
	\begin{align*}
		\cate{D}^\star(\mathcal{A}) & := \cate{K}^\star(\mathcal{A})\left[ (\mathrm{qis}^\star)^{-1} \right] \\
		& = \cate{K}^\star(\mathcal{A})/\mathcal{N}^\star .
	\end{align*}
	Kita menyebut $\cate{D}^+(\mathcal{A})$ (berturut-turut,
	$\cate{D}^-(\mathcal{A})$ dan $\cate{D}^\bdd(\mathcal{A})$) sebagai
	kategori turunan terbatas bawah (berturut-turut, terbatas atas dan terbatas)
	dari $\mathcal{A}$.
\end{definisi}

% segment-id: o014.li2.chapter4.derived-categories.p006
Persoalan yang sulit dihindari ialah mengendalikan ukuran kategori turunan
dalam arti teori himpunan, sebab kita tidak menginginkan
$\cate{D}(\mathcal{A})$ menjadi kategori ``besar''; lihat pembahasan terkait
dalam Catatan~\ref{rem:localization-size}. Persoalan ini ditangguhkan hingga
Korolari~\ref{prop:D-size}.
\index{masalah ukuran}

% segment-id: o014.li2.chapter4.derived-categories.pr003
\begin{proposisi}\label{prop:ses-vs-triangle}
	Setiap barisan eksak pendek
	$0 \to X \xrightarrow{f} Y \xrightarrow{g} Z \to 0$ dalam
	$\cate{C}(\mathcal{A})$ dapat diperluas secara kanonik menjadi segitiga
	terbedakan $X \to Y \to Z \xrightarrow{+1}$ dalam
	$\cate{D}(\mathcal{A})$. Secara tepat, dalam $\cate{K}(\mathcal{A})$
	terdapat diagram komutatif kanonik
% segment-id: o014.li2.chapter4.derived-categories.g006
	\[\begin{tikzcd}[column sep=large]
		X \arrow[r, "f"] \arrow[equal, d] & Y \arrow[equal, d] \arrow[r, "\alpha(f)"] & \Cone(f) \arrow[r, "\beta(f)"] \arrow[d, "\Phi"] & {X[1]} \arrow[dd, "{\Phi'}"] \\
		X \arrow[r, "f"] \arrow[d, "{\Phi'[-1]}"'] & Y \arrow[r, "g"] \arrow[equal, d] & Z \arrow[equal, d] & \\
		{\Cone(g)[-1]} \arrow[r, "{\beta(g)[-1]}"'] & Y \arrow[r, "g"'] & Z \arrow[r, "{-\alpha(g)}"'] & \Cone(g)
	\end{tikzcd}\]
	yang baris pertama dan ketiganya merupakan segitiga terbedakan, sedangkan
	$\Phi$ dan $\Phi'$ keduanya merupakan kuasi-isomorfisme.
\end{proposisi}

% segment-id: o014.li2.chapter4.derived-categories.pf006
\begin{bukti}
	Diagram komutatif tersebut hanyalah pernyataan ulang
	Lema~\ref{prop:cone-qis}. Baris pertama adalah segitiga terbedakan dari
	kerucut pemetaan $f$, sedangkan baris ketiga adalah rotasi kerucut pemetaan
	$g$.
\end{bukti}

% segment-id: o014.li2.chapter4.derived-categories.p007
Dengan demikian, barisan eksak pendek dalam kategori kompleks tercermin dalam
struktur bertriangulasi pada kategori turunan. Sebaliknya, kategori
abelian $\mathcal{A}$ yang membuat $\cate{D}(\mathcal{A})$ menjadi kategori
abelian sangatlah jarang; latihan-latihan bab ini akan memberikan
karakterisasinya.

% segment-id: o014.li2.chapter4.derived-categories.p008
Latihan-latihan bab ini juga akan memperkenalkan grup
$\mathrm{K}_0(\cate{D}^{\bdd}(\mathcal{A}))$ dan memberikan hubungan antara
grup tersebut dengan $\mathrm{K}_0(\mathcal{A})$ dari
Definisi~\ref{def:K0}; pembaca dianjurkan untuk mencobanya.

% segment-id: o014.li2.chapter4.derived-categories.p009
Berdasarkan karakterisasi $\mathcal{N}$ dan
Teorema~\ref{prop:Verdier-localization}, funktor kohomologis
$\Hm^n: \cate{K}(\mathcal{A}) \to \mathcal{A}$ mempunyai faktorisasi melalui
kategori turunan sebagai
% segment-id: o014.li2.chapter4.derived-categories.q013
\begin{equation*}
	\index[sym1]{Hn}
	\Hm^n: \cate{D}(\mathcal{A}) \to \mathcal{A}, \quad \Hm^n(X) = \Hm^0(X[n]),
\end{equation*}
dengan $n \in \Z$. Hal yang sama berlaku bagi
$\star \in \{+, -, \bdd\}$ dan $\cate{D}^{\star}(\mathcal{A})$.

% segment-id: o014.li2.chapter4.derived-categories.p010
Sifat universal menentukan funktor-funktor bertriangulasi
$\cate{D}^{\bdd}(\mathcal{A}) \to \cate{D}^{\pm}(\mathcal{A}) \to
\cate{D}(\mathcal{A})$ yang komutatif dengan funktor kohomologi $\Hm^0$.
Kita ingin memandang funktor-funktor tersebut sebagai pembenaman subkategori
bertriangulasi; untuk itu diperlukan hasil abstrak berikut.

% segment-id: o014.li2.chapter4.derived-categories.le001
\begin{lema}\label{prop:N-D-full-faithful}
	Misalkan $\mathcal{N}$ dan $\mathcal{I}$ subkategori bertriangulasi dari
	suatu kategori bertriangulasi $\mathcal{D}$, dengan $\mathcal{N}$ jenuh.
	Jika salah satu di antara syarat berikut berlaku dalam $\mathcal{D}$, maka
	funktor
	$i: \mathcal{I}/(\mathcal{N} \cap \mathcal{I}) \to \mathcal{D}/\mathcal{N}$
	penuh dan setia:
% segment-id: o014.li2.chapter4.derived-categories.l001
	\begin{enumerate}[(i)]
% segment-id: o014.li2.chapter4.derived-categories.i001
		\item setiap morfisme $Y \to N$ berfaktor sebagai $Y \to Y' \to N$;
% segment-id: o014.li2.chapter4.derived-categories.i002
		\item setiap morfisme $N \to Y$ berfaktor sebagai $N \to Y' \to Y$;
	\end{enumerate}
	di sini $Y \in \Obj(\mathcal{I})$ dan $N \in \Obj(\mathcal{N})$ diberikan,
	sedangkan $Y' \in \Obj(\mathcal{N} \cap \mathcal{I})$.
\end{lema}

% segment-id: o014.li2.chapter4.derived-categories.pf007
\begin{bukti}
	Berdasarkan dualitas, cukup ditinjau kasus ketika (i) berlaku. Ingat bahwa
	\eqref{eqn:S-N-D} menjamin
	$S(\mathcal{N} \cap \mathcal{I}) = S\mathcal{N} \cap \Mor(\mathcal{I})$.
	Kita hendak menerapkan kriteria dalam
	Proposisi~\ref{prop:subcat-localization} (ii). Karena itu, cukup dibuktikan
	pernyataan berikut: andaikan $s: W \to Y$ suatu morfisme
	dalam sistem multiplikatif $S\mathcal{N}$ dan
	$Y \in \Obj(\mathcal{I})$. Maka terdapat $g: V \to W$ sedemikian sehingga
	$V \in \Obj(\mathcal{I})$ dan $sg \in S\mathcal{N}$.

	Jelas terdapat segitiga terbedakan
	$W \xrightarrow{-s} Y \to N \xrightarrow{+1}$ dengan
	$N \in \Obj(\mathcal{N})$. Faktorkan $Y \to N$ sebagai
	$Y \xrightarrow{\alpha} Y' \xrightarrow{\beta} N$, dengan
	$Y' \in \Obj(\mathcal{N} \cap \mathcal{I})$. Perluas $\alpha$ menjadi
	segitiga terbedakan $V \to Y \xrightarrow{\alpha} Y' \xrightarrow{+1}$
	dalam $\mathcal{I}$ untuk memperoleh diagram berikut, yang bagian bergaris
	penuhnya telah ditentukan:
% segment-id: o014.li2.chapter4.derived-categories.g007
	\[\begin{tikzcd}
		Y \arrow[r, "\alpha"] \arrow[d, "\identity"] & Y' \arrow[d, "\beta"] \arrow[r] & TV \arrow[r] \arrow[dashed, d, "Tg"] & TY \arrow[d, dashed, "\identity"] \\
		Y \arrow[r] & N \arrow[r] & TW \arrow[r, "Ts"'] & TY 
	\end{tikzcd}\]
	dengan setiap baris merupakan segitiga terbedakan. Gunakan (TR4) untuk
	memilih $Tg: TV \to TW$ yang membuat seluruh diagram komutatif. Perhatikan
	bahwa $V \in \Obj(\mathcal{I})$ dan $TV \to TY$ termasuk dalam
	$S\mathcal{N}$; oleh karena itu, $sg: V \to Y$ juga termasuk dalam
	$S\mathcal{N}$. Pernyataan terbukti.
\end{bukti}

% segment-id: o014.li2.chapter4.derived-categories.pr004
\begin{proposisi}\label{prop:derived-full-faithful-subcat}
	Untuk semua $\star \in \{+, -, \bdd\}$, funktor
	$\cate{D}^{\star}(\mathcal{A}) \to \cate{D}(\mathcal{A})$ penuh dan setia.
	Untuk setiap $X \in \Obj(\cate{D}(\mathcal{A}))$, objek $X$ isomorfik dengan
	suatu objek dari $\cate{D}^{+}(\mathcal{A})$ (berturut-turut,
	$\cate{D}^{-}(\mathcal{A})$ dan $\cate{D}^{\bdd}(\mathcal{A})$) jika dan
	hanya jika $n \ll 0$ (berturut-turut, $n \gg 0$ dan $|n| \gg 0$)
	menyiratkan $\Hm^n(X) = 0$.
\end{proposisi}

% segment-id: o014.li2.chapter4.derived-categories.pf008
\begin{bukti}
	Bagian pertama merupakan penerapan Lema~\ref{prop:N-D-full-faithful}:
	di sana ambil $\mathcal{D} = \cate{K}(\mathcal{A})$,
	$\mathcal{I} = \cate{K}^\star(\mathcal{A})$, dan
	$\mathcal{N} = \mathcal{N}_{\Hm^0}$. Sebagai contoh, untuk $\star = +$,
		kita perlu menunjukkan bahwa bagi setiap
	$Y \in \Obj(\cate{K}^+(\mathcal{A}))$, $N \in \Obj(\mathcal{N})$, dan
	morfisme $N \to Y$, terdapat faktorisasi
% segment-id: o014.li2.chapter4.derived-categories.q014
	\[ N \to Y' \to Y, \quad Y' \in \Obj(\mathcal{N}^+). \]
	Hal ini berlaku karena, untuk $n \ll 0$, morfisme $N \to Y$ secara alami
	berfaktor sebagai
	$N \to \tau^{\geq n} N \to \tau^{\geq n} Y = Y$, dengan
	$\tau^{\geq n}$ funktor pemenggalan dari
	Definisi~\ref{def:truncation-cplx} dan
	$\tau^{\geq n} N \in \Obj(\mathcal{N}^+)$.

	Selanjutnya, jika $X \in \Obj(\cate{C}(\mathcal{A}))$ memenuhi
	$n \ll 0 \implies \Hm^n(X) = 0$, maka untuk $n \ll 0$ morfisme
	$X \to \tau^{\geq n} X$ merupakan kuasi-isomorfisme, sehingga merupakan
	isomorfisme dalam $\cate{D}(\mathcal{A})$. Ini mengarakterisasi, hingga
	isomorfisme, citra $\cate{D}^+(\mathcal{A}) \to \cate{D}(\mathcal{A})$.

	Argumen untuk $\star = -$ bersifat dual. Dengan cara yang sama,
	$\cate{K}^{\bdd}(\mathcal{A}) \to \cate{K}^\pm(\mathcal{A})$ menginduksi
	funktor penuh dan setia
	$\cate{D}^{\bdd}(\mathcal{A}) \to \cate{D}^\pm(\mathcal{A})$, dengan citra
	yang juga dikarakterisasi oleh $\Hm^n$. Dari sini diperoleh kasus
	$\cate{D}^{\bdd}(\mathcal{A}) \to \cate{D}(\mathcal{A})$.
\end{bukti}

% segment-id: o014.li2.chapter4.derived-categories.co003
\begin{korolari}
	Setiap funktor dalam rangkaian
	$\cate{D}^{\bdd}(\mathcal{A}) \to \cate{D}^{\pm}(\mathcal{A}) \to
	\cate{D}(\mathcal{A})$ merupakan pembenaman subkategori bertriangulasi.
\end{korolari}

% segment-id: o014.li2.chapter4.derived-categories.pf009
\begin{bukti}
	Gunakan barisan eksak panjang dari $\Hm^n$ untuk memverifikasi, bagi
	$\star \in \{+, -, \bdd\}$, bahwa $\cate{D}^\star(\mathcal{A})$ sebagai
	subkategori penuh aditif dari $\cate{D}(\mathcal{A})$ memenuhi syarat (ii)
	dalam Definisi--Proposisi~\ref{def:triangulated-subcat}. Dengan demikian,
	$\cate{D}^\star(\mathcal{A})$ menjadi subkategori bertriangulasi; kasus
	$\cate{D}^{\bdd}(\mathcal{A}) \to \cate{D}^{\pm}(\mathcal{A})$ juga jelas.
\end{bukti}

% segment-id: o014.li2.chapter4.derived-categories.d002
\begin{definisi}
	\index[sym1]{Dst@$\cate{D}^{[s,t]}(\mathcal{A})$, $\cate{D}^{\geq s}(\mathcal{A})$, $\cate{D}^{\leq t}(\mathcal{A})$}
	\index[sym1]{Kst@$\cate{K}^{[s,t]}(\mathcal{A})$, $\cate{K}^{\geq s}(\mathcal{A})$, $\cate{K}^{\leq t}(\mathcal{A})$}
	Untuk $-\infty \leq s \leq t \leq +\infty$, tuliskan
	$\cate{D}^{[s,t]}(\mathcal{A})$ (atau $\cate{K}^{[s, t]}(\mathcal{A})$)
	untuk subkategori penuh dalam $\cate{D}(\mathcal{A})$ (atau
	$\cate{K}(\mathcal{A})$) yang ditentukan oleh syarat
% segment-id: o014.li2.chapter4.derived-categories.q015
	\[ n \notin [s, t] \implies \Hm^n(X) = 0 . \]
	Selain itu, tetapkan
% segment-id: o014.li2.chapter4.derived-categories.q016
	\begin{align*}
		\cate{D}^{\geq s}(\mathcal{A}) & := \cate{D}^{[s, +\infty]}(\mathcal{A}), &
		\cate{D}^{\leq t}(\mathcal{A}) & := \cate{D}^{[-\infty, t]}(\mathcal{A}), \\
		\cate{K}^{\geq s}(\mathcal{A}) & := \cate{K}^{[s, +\infty]}(\mathcal{A}), &
		\cate{K}^{\leq t}(\mathcal{A}) & := \cate{K}^{[-\infty, t]}(\mathcal{A}).
	\end{align*}
\end{definisi}

% segment-id: o014.li2.chapter4.derived-categories.p011
Kita mempunyai funktor aditif alami
$\mathcal{A} \to \cate{C}(\mathcal{A}) \to \cate{D}(\mathcal{A})$ yang
memandang setiap objek sebagai kompleks yang terkonsentrasi pada derajat nol.
Teorema~\ref{prop:derived-cat-embedding} akan menunjukkan bahwa dengan
identifikasi ini $\mathcal{A}$ ekuivalen dengan
$\cate{D}^{\leq 0}(\mathcal{A}) \cap \cate{D}^{\geq 0}(\mathcal{A})$.

% segment-id: o014.li2.chapter4.derived-categories.r001
\begin{catatan}
	Jika $\mathcal{A}$ bersifat $\Bbbk$-linear, dengan $\Bbbk$ suatu gelanggang
	komutatif, maka $\cate{D}(\mathcal{A})$, $\cate{D}^{\bdd}(\mathcal{A})$,
	dan seterusnya, serta pembenaman-pembenaman di antaranya, juga bersifat
	$\Bbbk$-linear; lihat Teorema~\ref{prop:localization-additivity} (iii).
\end{catatan}

% segment-id: o014.li2.chapter4.derived-categories.p012
Karena funktor pemenggalan $\tau^{\leq n}$ dan $\tau^{\geq n}$ dari
Definisi--Proposisi~\ref{def:K-truncation-functors} mempertahankan kompleks
asiklik, keduanya menginduksi
$\tau^{\leq n}: \cate{D}(\mathcal{A}) \to \cate{D}^{\leq n}(\mathcal{A})$
dan
$\tau^{\geq n}: \cate{D}(\mathcal{A}) \to \cate{D}^{\geq n}(\mathcal{A})$.

% segment-id: o014.li2.chapter4.derived-categories.pr005
\begin{proposisi}\label{prop:derived-truncation-adjunction}
	Untuk setiap $n \in \Z$, kita mempunyai pasangan-pasangan adjoin kanonik
	dan segitiga terbedakan
% segment-id: o014.li2.chapter4.derived-categories.g008
	\begin{equation*}\begin{gathered}\begin{tikzcd}
		\text{inklusi}: \cate{D}^{\leq n}(\mathcal{A}) \arrow[shift left, r] & \cate{D}(\mathcal{A}) :\tau^{\leq n} \arrow[shift left, l] ,
	\end{tikzcd} \\
	\begin{tikzcd}
		\tau^{\geq n}: \cate{D}(\mathcal{A}) \arrow[shift left, r] & \cate{D}^{\geq n}(\mathcal{A}) :\text{inklusi} \arrow[shift left, l] ,
	\end{tikzcd} \\
		\tau^{\leq n} X \to X \to \tau^{\geq n+1} X \xrightarrow{+1}, \quad X \in \Obj(\cate{D}(\mathcal{A})).
	\end{gathered}\end{equation*}
\end{proposisi}

% segment-id: o014.li2.chapter4.derived-categories.pf010
\begin{bukti}
	Cukup dibahas pasangan adjoin pertama. Misalkan
	$X \in \Obj(\cate{D}^{\leq n}(\mathcal{A}))$ dan
	$Y \in \Obj(\cate{D}(\mathcal{A}))$. Unit dan kounit dari pasangan adjoin
	yang dimaksud berturut-turut adalah
	$\eta_X: X \to \tau^{\leq n} X$ dan
	$\varepsilon_Y: \tau^{\leq n} Y \to Y$. Morfisme $\varepsilon_Y$ adalah
	morfisme kanonik yang sudah ada pada tingkat kompleks, sedangkan $\eta_X$
	adalah invers dari kuasi-isomorfisme $\tau^{\leq n} X \to X$. Tinggal
	memverifikasi identitas-identitas segitiga. Setelah melakukan pemenggalan
	pada tingkat kompleks, cukup meninjaunya ketika $X$ berasal dari
	$\cate{C}^{\leq n}(\mathcal{A})$; dalam kasus ini
	$\eta_X = \identity_X$. Identitas-identitas segitiga yang diperlukan
	kemudian mengikuti dari versi kompleks dalam
	Proposisi~\ref{prop:truncation-adjunction}.

	Adapun mengenai segitiga terbedakan tersebut, ingat bahwa pada tingkat
	$\cate{C}(\mathcal{A})$ terdapat barisan eksak pendek
	$0 \to \tau^{\leq n} X \to X \to \tilde{\tau}^{\geq n+1} X \to 0$ dan
	kuasi-isomorfisme
	$\tilde{\tau}^{\geq n+1} X \to \tau^{\geq n+1} X$
	(Lema~\ref{prop:truncation-ses}).
\end{bukti}

% segment-id: o014.li2.chapter4.derived-categories.r002
\begin{catatan}[Konstruksi Langsung Kategori Turunan]
	Nyatakan dengan $\mathrm{Qis} \subset \Mor(\cate{C}(\mathcal{A}))$ himpunan
	semua kuasi-isomorfisme. Andaikan dalam definisi $\cate{D}(\mathcal{A})$
	kita tidak terlebih dahulu mengambil $\cate{K}(\mathcal{A})$, melainkan
	langsung menambahkan invers bagi kuasi-isomorfisme untuk memperoleh
	$\cate{C}(\mathcal{A})[\mathrm{Qis}^{-1}]$. Apakah hasilnya ekuivalen dengan
	$\cate{D}(\mathcal{A})$? Jawabannya ya. Dengan sedikit memikirkan sifat
	universal, terlihat bahwa kuncinya ialah memverifikasi sifat berikut bagi
	setiap kategori $\mathcal{D}$ dan funktor
	$G: \cate{C}(\mathcal{A}) \to \mathcal{D}$: andaikan $G$ memetakan
	$\mathrm{Qis}$ ke isomorfisme, sedangkan
	$f, g \in \Hom_{\cate{C}(\mathcal{A})}(X, Y)$ saling homotopik. Maka
	$Gf = Gg$.

	Kuncinya ialah menggunakan silinder pemetaan. Misalkan $h$ suatu homotopi
	dari $f$ ke $g$. Dengan cara dalam
	Proposisi~\ref{prop:cylinder-homotopy} (ii), homotopi tersebut menentukan
	morfisme $\tilde{h}: \Cyl_X \to Y$ sedemikian sehingga
	$f = \tilde{h} i_0$ dan $g = \tilde{h} i_1$. Akan tetapi, bagian (i) dari
	proposisi tersebut menyatakan $j i_0 = j i_1$, dengan
	$j: \Cyl_X \to X$ suatu kuasi-isomorfisme. Karena itu, $Gj$ merupakan
	isomorfisme, dan hal ini pada gilirannya mengakibatkan
	$G i_0 = G i_1$. Sifat tersebut terbukti.

	Meskipun $\cate{D}(\mathcal{A})$ dapat dibangun langsung sebagai
	$\cate{C}(\mathcal{A})[\mathrm{Qis}^{-1}]$, keuntungan melewati
	$\cate{K}(\mathcal{A})$ ialah bahwa kategori terakhir mempunyai struktur
	bertriangulasi; alasan lainnya ialah bahwa $\mathrm{Qis}$ bukan suatu sistem
	multiplikatif. Pernyataan yang sama berlaku bagi
	$\cate{D}^\star(\mathcal{A})$, dengan $\star \in \{+, -, \bdd\}$.
\end{catatan}

% segment-id: o014.li2.chapter4.derived-categories.p013
Jenis penting lainnya dari subkategori bertriangulasi ialah subkategori yang
ditentukan oleh kohomologi.

% segment-id: o014.li2.chapter4.derived-categories.d003
\begin{definisi}[Subkategori Bertriangulasi yang Ditentukan oleh Kohomologi]
	\index[sym1]{DTA@$\cate{D}_{\mathcal{T}}(\mathcal{A})$}
	Misalkan $\mathcal{T}$ suatu subkategori Serre lemah dari $\mathcal{A}$
	(Definisi~\ref{def:Serre-subcat}). Definisikan subkategori penuh
	$\cate{D}_{\mathcal{T}}(\mathcal{A})$ dari $\cate{D}(\mathcal{A})$ dengan
% segment-id: o014.li2.chapter4.derived-categories.q017
	\[ X \in \Obj(\cate{D}_{\mathcal{T}}(\mathcal{A})) \iff \forall n \in \Z, \; \Hm^n(X) \in \Obj(\mathcal{T}). \]
	Menurut Proposisi~\ref{prop:cut-subcat-by-cohomologies}, ini merupakan
	subkategori bertriangulasi jenuh dari $\cate{D}(\mathcal{A})$.

	Untuk $\star \in \{+, -, \bdd\}$, definisikan subkategori bertriangulasi
	jenuh
	$\cate{D}^\star_{\mathcal{T}}(\mathcal{A}) := \cate{D}_{\mathcal{T}}(\mathcal{A}) \cap \cate{D}^\star(\mathcal{A})$.
\end{definisi}

% segment-id: o014.li2.chapter4.derived-categories.p014
Terakhir, terdapat hubungan sederhana antara kategori turunan
$\mathcal{A}$ dan $\mathcal{A}^{\opp}$. Berikan struktur bertriangulasi pada
kategori lawan dari suatu kategori bertriangulasi dengan cara dalam
Catatan~\ref{rem:triangulated-op}.

% segment-id: o014.li2.chapter4.derived-categories.pr006
\begin{proposisi}\label{prop:derived-cat-op}
	Ekuivalensi
	$\sigma: \cate{C}(\mathcal{A}^{\opp}) \rightiso \cate{C}(\mathcal{A})^{\opp}$
	dari Definisi--Proposisi~\ref{def:sigma} menginduksi ekuivalensi
	kategori-kategori bertriangulasi
% segment-id: o014.li2.chapter4.derived-categories.q018
	\[ \cate{D}(\mathcal{A}^{\opp}) \simeq \cate{D}(\mathcal{A})^{\opp}, \quad \cate{D}^\pm(\mathcal{A}^{\opp}) \simeq \cate{D}^\mp(\mathcal{A})^{\opp}, \quad \cate{D}^{\bdd}(\mathcal{A}^{\opp}) \simeq \cate{D}^{\bdd}(\mathcal{A})^{\opp}. \]
	Selain itu, untuk semua $n \in \Z$ berlaku
	$\tau^{\leq n} \circ \sigma = \sigma \circ \tau^{\geq -n}$ dan
	$\tau^{\geq n} \circ \sigma = \sigma \circ \tau^{\leq -n}$.
\end{proposisi}

% segment-id: o014.li2.chapter4.derived-categories.pf011
\begin{bukti}
	Proposisi~\ref{prop:sigma-homotopy} menunjukkan bahwa $\sigma$ menginduksi
	$\cate{K}(\mathcal{A}^{\opp}) \simeq \cate{K}(\mathcal{A})^{\opp}$.
	Funktor tersebut mempertahankan segitiga terbedakan (bandingkan dengan
	cermat diagram dalam Proposisi~\ref{prop:sigma-triangulated} dan
	Catatan~\ref{rem:triangulated-op}: setelah baris pertama pada diagram yang
	disebut lebih dahulu dibalik, baris itu merupakan citra segitiga terbedakan;
	baris kedua merupakan segitiga terbedakan dalam $\cate{K}(\mathcal{A})$ dan,
	setelah dibalik, menjadi segitiga terbedakan dalam
	$\cate{K}(\mathcal{A})^{\opp}$). Funktor tersebut juga mempertahankan
	kohomologi (Catatan~\ref{rem:sigma-vs-cohomology}), sehingga menginduksi
	$\cate{D}(\mathcal{A}^{\opp}) \rightiso \cate{D}(\mathcal{A})^{\opp}$.
	Pernyataan mengenai funktor pemenggalan mengikuti dari
	Catatan~\ref{rem:truncation-sigma}.
\end{bukti}
